\section{はじめに}
\label{sec:introduction}

新しいコーパスに配備した感情認識器は，学習データに含まれなかった話者，録音条件，感情の引き出し方，注釈の慣習に出会い，コーパスをまたぐと正解率が下がる\cite{schuller2010crosscorpus,gideon2021addog,feng2023fewshot}．
新しいコーパスにラベルを付けるには費用がかかるので，実務家は通常，感情ごとに少数の発話にラベルを付け，他のコーパスのラベル付きの発話を加える．
それらのソースコーパスが役立つか，どれだけ役立つかは，コーパスとターゲットのラベルの本数によって決まる．
マルチモーダル感情認識では，ターゲットの発話が複数のモダリティを伴うことが多く，コーパスをまたいで学習する学習器はこれを第2の情報源として使える．

この設定を扱う研究には2つの系統があり，それぞれ別の面から取り組んでいる．
複数のドメインからの学習の理論は，ソースを合併するとターゲットのラベルが少ないときは役立ち，多いときは害になることを示し，ターゲットの誤差を，ソースの誤差，分布の隔たり（divergence），ターゲットの標本の大きさで上から抑える\cite{bendavid2010theory,blitzer2007learning,desilva2023value}．
少数のターゲットラベルによる感情認識の手法は，音声のモデルや音声とテキストのモデルを少数のラベル付き発話で適応させ\cite{feng2023fewshot,he2025m4ser,zaidi2024enhancing}，半教師あり学習と蒸留は，ラベルのないデータやより強い先生（teacher）を使う\cite{lee2013pseudolabel,hinton2015distilling,lopezpaz2016unifying}．
少数のターゲットラベルによるコーパス横断の設定において，どちらの系統も，少数のターゲットラベルが分類器のどの部分を直すか，ソースがラベルの代わりにどれだけを補うか，配備したモデルが1つのモダリティしか見ないときに第2のモダリティが何を加えうるかを示していない．
顔を含む複数のモダリティを使うコーパス横断の手法も，依然として少ない\cite{ahn2025multitask}．

本論文は，凍結した（微調整しない）エンコーダの特徴の上の線形ヘッドについて，3つの問いに答える．
この設定は分析できるほど単純であり，4つのターゲットで音声エンコーダを低ランクの適応で微調整した確認では，微調整より強い．
(Q1)~少数のターゲットラベルは，分類器の何を直すか．
(Q2)~ソースは，ラベルの代わりにどれだけを補うか．
(Q3)~その答えから，同じ手順で調整した標準的な半教師あり学習と蒸留の手法を上回る手法が得られるか．
6つのコーパスから作った7つの群を，音声・顔・テキストの特徴とそれらの融合で調べる（図\ref{fig:overview}）．
4クラス分類器の感情の対ごとの判定を判断基準のずれ（threshold offset）と識別力（separation）で表し，それぞれについて感情ごとに$K$本のラベルでのコストを導き，その予言を，コーパスのすべての順序対の利得の地図（gain map）と照らし合わせる．

\begin{figure*}[t]
\centering
\FIG{overview}{\textwidth}
\caption{概観．出発点の分類器は，凍結したエンコーダを用いて，ラベル付きのソースコーパスで学習する．第I部：単一モダリティの音声と顔の特徴で$K \le 5$のとき，感情ごとに$K$本のラベルが加える改善の大半は判断基準から生じ（残る誤差は$1/(2K)$），音声の向きはゆっくり改善する（残る誤差は$q/K$）．第II部：ソースだけで学習し，生徒がもたないモダリティを含む融合の分類器が，ラベルのないターゲットの発話すべてに，感情ごとの合計の重みが等しい弱いラベルを付け，出発点にアンカーした線形の生徒が，この弱いラベルと$K$本のラベルから学ぶ．テストのとき，生徒は自分のモダリティだけを必要とする．}
\label{fig:overview}
\end{figure*}

本論文の貢献は次のとおりである．
\begin{enumerate}
\item 顔を含む3つのモダリティとその融合について，7つの群のすべての順序対と，感情ごとに1本からプール全体までのラベルの本数にわたって作った，コーパス横断の合併の利得の地図．30回の反復にもとづく対応のある区間を含む（\ref{sec:part1-map}節）．
\item 利得の判断基準の部分と向きの部分への分解．分解は$K$の関数として測り，同じ分割の上でラベルを使わない中心化と比べる．単一モダリティの音声と顔で$K \le 5$のとき，ラベルが合併の分類器に加える改善の\NUM{p1-e29-critshare-kle5-single-lo}--\NUM{p1-e29-critshare-kle5-single-hi}が判断基準から生じ，音声の分類器の向きは$q/K$の速さでしか改善しない（測った$q$は\NUM{p1-qdir-audio-lo}--\NUM{p1-qdir-audio-hi}）（\ref{sec:part1-fix}節）．
\item ソースの価値をターゲットのラベルの本数に換算する式$n = q e^2/(1-e^2)$．この式は推定した向きと標本サイズの古典的な関係を逆に解いて得たものであり，測った値は，音声のソース1つで感情ごとに\NUM{p1-e32-n1src-audio-lo}--\NUM{p1-e32-n1src-audio-hi}本，音声のソースすべてを合併すると\NUM{p1-e32-nall-audio-lo}--\NUM{p1-e32-nall-audio-hi}本である（\ref{sec:part1-value}節）．
\item ソースだけで学習し，ターゲットのデータで学び直さず，生徒がもたないモダリティを見る融合の分類器から，クラスの釣り合いをとった弱いラベルをラベルのないターゲットのプール全体に付け，仮のターゲットで$K$ごとに選んだハイパーパラメータで，アンカーつきの線形の生徒を学習する手法．基本の手法は，生徒と$K \in \{0, 5\}$の\NUM{cmp-ours-vs-best-cells-n}の組み合わせのうち\NUM{cmp-ours-vs-best-pos-n}で，同じ手順で調整した最も強い比較手法を上回る．ターゲットのプールで学び直した先生を加え，弱いラベルの重みづけを$K$ごとに選ぶ拡張版は，そのすべての組み合わせで最も強い比較手法を\NUM{cmp-oursx-vs-best-range-lo}--\NUM{cmp-oursx-vs-best-range-hi}ポイント上回り，第2のエンコーダでも上回る．本論文で測った場合には，先生の出発点に対するUARの優位が利得を上から抑える（\ref{sec:part2-crossmodal}節と\ref{sec:comparison}節）．
\end{enumerate}
我々の知る限り，これら4つの設定と測定の組み合わせはこれまで報告されていない．
\ref{sec:related}節では，各構成要素と各式を先行研究に対して位置づける．

\ref{sec:related}節で関連研究を概観し，\ref{sec:setting}節で手順を定義し，\ref{sec:part1-what-few-labels-fix}節で第I部，\ref{sec:part2-mechanism-to-method}節で第II部，\ref{sec:comparison}節で比較を示す．
\ref{sec:discussion}節と\ref{sec:limitations-ethics}節では，結果の意味と限界を論じる．
